\chapter{宏观潜空间：任务压缩、受限几何与闭环控制}
\label{chp:macro_latent_space_topology}

一个智能系统若要在有限带宽上维持跨任务、跨时间的控制，就不能把全部局部活动逐项送入宏观回路；它必须形成某种低维状态，用来概括当前目标、风险、资源、承诺与环境变化。然而，低维并不意味着存在一个脱离内容、永远稳定的“纯几何自我”，球面、李群或拓扑商也不能仅凭形状相似便成为智能的必然本体。本章把历史论述中的重整化、超球面、非对易、切空间、绝热分离、自指和向下作用依次改写为可检验的计算模型。我们先定义压缩相对于哪些查询有效，再区分方向变量、幅度变量和资源账本；随后给出约束流形上的连续与离散更新、快慢回路的成立条件、闭环不动点与极限环的不同判据，以及宏观状态如何经接口改变底层增益、约束和探索策略。HSF-HD在这里提供的是智能一般状态动力学的一组模型族；AgentOS的SPC、TCE、CCM、Boundary和外部记忆只是其中一种工程实例，而不是所有智能系统必须采用的唯一结构。

\section{任务商空间：从高维活动到可验证宏观状态}
\label{sec:meta_topology_rg_flow}

我们先把“宏观潜空间”定义成任务相关的状态压缩，而不是从物理重整化群直接移植出来的实体。设时刻$t$的联合微观状态为$x_t=(h_t,G_t,m_t,e_t)\in\mathcal X$，其中$h_t\in\mathbb R^{d_h}$是快速分布表示，$G_t$是带身份、关系、来源和版本的局部结构网络，$m_t$是内部或外部记忆记录，$e_t$是可观测环境上下文。四类对象分别遵守连续向量、离散图、版本化记录和外部过程的更新规则，不能因为都被写入一个状态元组就假定它们具有同一线性结构。对任务族$\mathcal T$，我们给定查询集合$\mathcal Q_{\mathcal T}$、允许动作集合$\mathcal A_{\mathcal T}$和风险度量$d_q$。两个状态$x,x'\in\mathcal X$只有在所有声明查询与干预上都近似不可区分时才视为同一宏观类：
\begin{equation}
 x\sim_{\mathcal T,\varepsilon}x'
 \Longleftrightarrow
 \sup_{q\in\mathcal Q_{\mathcal T}}d_q\!\left(q(x),q(x')\right)\leq\varepsilon_q,
 \quad
 \sup_{a\in\mathcal A_{\mathcal T}}d_a\!\left(P_x^a,P_{x'}^a\right)\leq\varepsilon_a .
\end{equation}
这里$P_x^a$表示从状态$x$执行动作$a$后的结果分布。第一项检查预测查询，第二项检查干预后果；若只比较静态读出，两个具有相似表面描述但因果响应不同的状态会被错误合并。于是，形式上的商空间$\mathcal Z_{\mathcal T}=\mathcal X/\!\sim_{\mathcal T,\varepsilon}$只有在近似关系足够稳定、类内误差可控并且任务族固定的区间内才有操作意义。任务、观测或动作集合改变时，等价类也会改变，不存在一次压缩后对所有未来问题都无损的绝对商空间。

工程实现不需要显式枚举等价类。我们使用编码器$R_{\theta,t}:\mathcal X\rightarrow\mathbb R^D$得到宏观状态$z_t=R_{\theta,t}(x_t)$，并用一组任务读出器$D_{q,t}$复原需要的量。编码参数$\theta$、结构版本$G_t$、输入分布和度量都可能随时间变化，因此压缩质量必须写成带版本的风险：
\begin{equation}
 \mathcal L_{\mathrm{macro}}(t)=
 \mathbb E_{q\sim\mu_t}\!\left[d_q\!\left(D_{q,t}(z_t),q(x_t)\right)\right]
 +\lambda_{\mathrm{int}}\mathcal L_{\mathrm{intervention}}(t)
 +\lambda_{\mathrm{cal}}\mathcal L_{\mathrm{calibration}}(t).
\end{equation}
权重只有在三项经过协议换算或均被定义为无量纲风险时才可直接相加；否则应采用硬约束、词典序或Pareto比较。所谓“宏观序参量”因此是经验证的摘要变量：它能在规定预算内支持某些查询、规划与控制，但不等于全局波函数的拓扑不变量积分。局部结构网络与全局分布表示在这里分工明确：前者保存对象身份、路径和约束，后者表示候选、活动与不确定性；编码器必须通过绑定表将两者对应起来，并在来源冲突或结构版本变化时重新校准。

这种定义保留了重整化直觉中有价值的一点：宏观控制不必跟随每个高频自由度。但我们不把频率滤除本身当作意义提取。一个持续时间很短的安全告警可能对任务极其关键，一个长期重复的背景信号也可能完全无关。压缩门控应依据任务价值、预测增益、干预影响和不可逆风险，而不是仅依据“高频”或“低频”。验证时至少要做四类试验：同类状态的压缩一致性、不同因果状态的可分辨性、分布漂移后的校准、以及删除某类输入后的反事实退化。只有这些试验通过，我们才说宏观潜空间在相应任务域内构成有效粗粒化；它仍然是一种模型假设和工程结果，不能由降维成功推出系统与世界整体同构。

压缩还要面对任务之间的迁移。若两个任务共享对象身份却使用不同风险函数，可以复用结构绑定而重新训练读出；若它们连可观测变量和动作语义都不同，则复用同一潜空间可能产生负迁移。工程上应保存任务描述、训练分布、适用查询和失败样本，利用跨任务探针判断哪些坐标真正可复用。新任务到来时，先比较旧表示的查询误差和干预误差，再决定扩充维度、建立新分支或仅增加适配器。宏观空间的稳定来自受控迁移与版本兼容，而不是把一个低维编码宣布为永恒不变量。

\section{受限几何：方向、幅度与资源状态的分账}
\label{sec:hypersphere_geometry_constraint}

球面约束可以成为某些宏观变量的有用设计，却不是智能稳定性的必然定理。设宏观表示被分解为方向变量$u_t\in\mathbb R^D$、置信或重要性幅度$r_t\in\mathbb R_{\geq0}$、资源状态$b_t\in\mathbb R^k$与离散模式$c_t\in\mathcal C$。若任务只关心相对偏好方向，我们可以令$\|u_t\|_2=1$，使$u_t\in\mathbb S^{D-1}$；但$r_t$仍记录证据强度或控制置信，$b_t$记录时延、显存、调用配额与设备功率，$c_t$记录安全模式或承诺阶段。完整宏观状态因而更接近乘积空间
\begin{equation}
 \mathcal Z=\mathbb S^{D-1}\times\mathbb R_{\geq0}\times\mathcal B\times\mathcal C,
 \qquad z_t=(u_t,r_t,b_t,c_t),
\end{equation}
而不是一个能够吞并全部变量的单位球。方向归一化消除了共同尺度自由度，却同时丢失了幅度信息；如果“轻微怀疑”和“高度确信”被归一化为同一方向，控制器必须从$r_t$或其他通道恢复差别。把所有幅度剥离会造成欠表达，而非天然安全。

采用球面需要明确三个条件。第一，方向确实具有尺度不变性，即正比例缩放不改变相关决策；第二，零向量与低信噪比状态另有拒绝表示，否则归一化会把噪声放大成确定方向；第三，所用内积与特征单位经过定义。若坐标来自不同量纲或不同校准区间，直接使用欧氏范数没有意义。可引入正定度量$M_t\succ0$并定义$\|u\|_{M_t}^2=u^\top M_tu=1$，但当$M_t$变化时，约束的总导数为
\begin{equation}
 \frac{d}{dt}(u_t^\top M_tu_t)=2u_t^\top M_t\dot u_t+u_t^\top\dot M_tu_t=0.
\end{equation}
因此，仅令更新方向与$u_t$在当前欧氏内积下正交并不足以维持时变度量中的单位范数。结构、任务和校准协议改变引起的$\dot M_t$必须进入更新，或者在离散提交时重新投影并记录产生的语义失真。

形与质在本书中可继续作为关系结构与属性内容的分析语言，但二者并非绝对正交。局部结构网络$G_t$的边类型会改变全局分布$h_t$的传播，而活动证据又会触发结构边的新增、降权或撤销；宏观方向$u_t$与幅度$r_t$也存在可识别耦合。我们可以通过门控函数$g(r_t,b_t,c_t)$限制方向更新的增益，或在证据不足时维持原状态：
\begin{equation}
 \dot u_t=g(r_t,b_t,c_t)\,v_t,
 \qquad v_t\in T_{u_t}\mathbb S^{D-1}.
\end{equation}
这里$g$是无量纲增益，$v_t$的单位为每秒，二者都不是焦耳能量。设备实际能耗由功率计量或硬件模型给出，不能从$\|u_t\|$守恒推出“不创造也不毁灭能量”。同样，范数不变只排除了径向发散，并不排除沿球面的高速振荡、错误吸引子、语义漂移或危险输出。

反例能显示边界。一个归一化控制向量可以在两个相反策略之间快速翻转，范数始终为一，系统却完全不稳定；一个允许幅度缓慢变化的控制器可以借助置信衰减安全退出，反而更鲁棒。若宏观任务需要表示“无偏好”“多目标冲突”或“停止控制”，球面上的单一单位向量也不够，必须增加拒绝状态、混合分布或集合值表示。因此，本章把超球面视为可选约束族：当尺度不变性、非零信号、固定度量和适当读出均得到验证时，它可降低参数冗余；在其他情形中，欧氏空间、概率单纯形、双曲空间或离散状态机都可能更合适。几何由任务与接口选择，而不是先验宣布哪一种几何等于“自我”。

即使球面模型适用，也要说明坐标变换下哪些读出保持不变。若任意正交变换都会同步改变解码器，则单个坐标轴没有固有语义；解释必须落在可验证的子空间、方向关系或任务读出上。我们可通过旋转探针、坐标置换和噪声扰动检验表示是否依赖偶然基底，并在接口层保存坐标版本。不同模型之间的方向相似只有经过对齐映射和置信区间后才可比较，不能把余弦相近直接说成价值相同。这样，几何约束才从视觉上漂亮的图形变成可迁移、可审计的表示契约。

\section{非交换更新：历史顺序、算子组合与经验保持}
\label{sec:non_abelian_lie_group_properties}

认知历史具有顺序效应，但顺序效应不自动要求非阿贝尔李群。我们先在一般状态空间上定义两个合法更新算子$F_A$与$F_B$，分别代表事件$A$和$B$在给定上下文、权限和记忆版本下的处理。若$F_B(F_A(z))\neq F_A(F_B(z))$，则系统对这两个事件存在非交换更新。造成差异的机制可能是参数学习、缓存覆盖、不可逆外部行动、结构图编辑或资源消耗；其中许多算子既不可逆也不光滑，不能纳入群。李群模型只适用于一族可逆、光滑且具有共同单位元的变换。把一切路径依赖都称为群旋转，会隐藏删除、提交和失败补偿等真实工程语义。

在确实存在连续对称性的局部模式内，我们可令$U_A(s)=\exp(sX_A)$、$U_B(s)=\exp(sX_B)$，其中$X_A,X_B$属于李代数$\mathfrak g$，参数$s$为无量纲更新强度。小步长下，Baker--Campbell--Hausdorff展开给出
\begin{equation}
 \log\!\left(e^{\delta X_B}e^{\delta X_A}\right)
 =\delta(X_A+X_B)+\frac{\delta^2}{2}[X_B,X_A]+O(\delta^3).
\end{equation}
对易子$[X_B,X_A]$刻画顺序差异的二阶局部贡献；只有当生成元、坐标图和误差阶数都已声明时，这个量才可解释。若事件包含离散图编辑，正确模型是带版本的算子序列，例如$G_{k+1}=E_B^{v+1}(E_A^v(G_k))$，并比较保留查询、冲突集合和回滚结果，而不是强行寻找连续生成元。

创伤后再安抚与安抚后再遭遇创伤可能留下不同状态，这一例证应被理解为历史依赖假说。要把它变成经验主张，需要固定观测窗口、输入强度、初始状态与测量指标，区分记忆写入、价值更新、生理调节和外部环境改变的贡献。不同结果并不证明存在某个特定球面，也不证明人格必然由几何相位构成。若连续路径$\gamma$位于具有联络$A_t$的向量丛上，可用路径有序指数
\begin{equation}
 U_\gamma=\mathcal P\exp\!\left(\int_0^T A_t\,dt\right)
\end{equation}
描述沿途变换；闭路后的$U_\gamma$可作为回到相似观测状态后内部表示仍发生变化的诊断。但该和乐依赖所选联络、规范和可观测读出，不等同于智慧、人格或真理。一个系统也可能因为缓存污染形成巨大路径效应，这显然不能称作阅历增长。

对工程而言，历史保持需要比非交换性更多的条件。每次宏观更新应记录事件来源、时间、状态版本、应用顺序、执行结果和可撤销性；需要交换的并发操作应通过幂等键、交换律测试或CRDT式规则减少顺序敏感；不可交换操作则必须串行化、建立前置条件或提供补偿事务。局部结构网络保存“发生了什么、作用于谁以及依赖哪一版本”，全局分布表示保存“哪些候选被激活以及不确定性如何变化”。两者经身份绑定后，才能区分真正的历史学习与单纯的活动残留。历史记录还应区分原始事件、系统解释与后续评价，避免模型把自己生成的解释再次当作独立事实。

李群表示也有范围限制。旋转算子适合表达保范数的坐标变换，却不能表达删除节点、生成新身份、工具调用失败或外部承诺生效。若一段历史同时含连续适应和离散事件，可使用混合自动机：连续模式内用流形动力学，事件边界用带守卫条件的跳变映射。跳变前后要验证身份映射和关键约束，而不能假定两个坐标图自然同构。由此，非交换性的真正信息在于“哪些操作次序不能互换以及为什么”，而不在于为所有经验附加一个量子或几何名称。

验证协议包括：交换次序的配对实验、不同步长下对易子估计、离散编辑的重放一致性、版本冲突注入、以及删除某段历史后的反事实读出。如果顺序效应随步长消失，它可能只是数值误差；如果效应只来自外部不可逆操作，就应归入环境状态；如果多次重放无法复现，则需要检查随机种子与未记录输入。只有能被定位、复现并与任务结果关联的路径依赖，才构成宏观潜空间模型的有效组成，而不是由漂亮的群论形式替代因果分析。

\section{切空间动力学：约束梯度、离散回缩与安全边界}
\label{sec:tangent_space_dynamics}

设方向状态$u(t)\in\mathbb S^{D-1}$，底层经任务压缩得到的驱动为$s(t)=R_{\theta,t}(x_t)\in\mathbb R^D$。若度量固定为欧氏内积，切空间投影为$P_u=I-uu^\top$，连续模型可写成
\begin{equation}
 \dot u(t)=\eta(t)P_{u(t)}s(t),
 \qquad \|u(0)\|_2=1,
\end{equation}
其中$\eta(t)$的单位为驱动量倒数每秒，使$\dot u$具有每秒单位。由$u^\top P_us=0$可得$d\|u\|_2^2/dt=0$。这是一个明确的数学结果：在解存在且数值积分精确的连续模型里，投影流保持范数。它不证明任务目标下降，更不证明控制安全。$s(t)$可能指向错误方向，$\eta$可能过大，沿面速度可能无界，观测延迟也可能造成持续振荡。

若更新来自可微任务函数$J(u,t)$，我们使用黎曼梯度$\operatorname{grad}J=P_u\nabla_uJ$。当$J$还依赖目标$y_t$、结构$G_t$、输入分布$\mu_t$与度量$M_t$时，沿轨迹的总导数必须包含这些变化：
\begin{equation}
 \frac{dJ}{dt}=\langle\nabla_uJ,\dot u\rangle
 +\left\langle\partial_GJ,\dot G\right\rangle
 +\left\langle\partial_yJ,\dot y\right\rangle
 +\left\langle\partial_\mu J,\dot\mu\right\rangle
 +\partial_tJ.
\end{equation}
在固定目标、结构、输入分布且$\dot u=-\eta\operatorname{grad}J$时，第一项非正；其他项存在时，局部梯度下降未必带来总目标下降。离散图$G$没有普通导数，实际实现应在事件时刻计算跳变量$J(u^+,G^+)-J(u^-,G^-)$，并通过影子版本、保留测试和回滚门槛控制结构提交。

连续保范数也不会自动传给欧拉离散。更新$\tilde u_{k+1}=u_k+\Delta t\,\eta_kP_{u_k}s_k$一般满足$\|\tilde u_{k+1}\|>1$。可采用回缩
\begin{equation}
 u_{k+1}=\operatorname{Retr}_{u_k}(v_k)
 =\frac{u_k+v_k}{\|u_k+v_k\|_2},
 \qquad v_k=\Delta t\,\eta_kP_{u_k}s_k,
\end{equation}
并要求$u_k+v_k\neq0$。步长应依据驱动上界、局部Lipschitz常数、延迟和容许角位移选择；在异步系统中还要记录$s_k$对应的状态版本，拒绝过期驱动。若采用指数映射，计算更符合几何但代价更高；若采用简单归一化，必须测量回缩引入的读出误差。任何实现都要提供拒绝选项，不能在驱动置信不足时仍强迫单位向量选择一个方向。

历史公式中的反对称生成元可以在实向量情形下重新解释。令$\Omega=s u^\top-u s^\top$，则$\Omega^\top=-\Omega$且$\Omega u=s-(u^\top s)u=P_us$。因此$\dot u=\eta\Omega u$与投影方程等价。这个构造说明保范数流可写为瞬时旋转，却不能推出“底层能量只会让自我转得更快”。$\|s\|$是编码驱动的数值尺度，不是焦耳能量；过大的$\|s\|$会导致离散积分失稳、控制饱和和接口拥塞。CCM式监控应限制角速度、创新率、过期比例和资源占用，Boundary应限制宏观命令的权限与作用域，TCE式控制器还需根据风险降低增益或进入安全模式。

约束几何还需处理多个候选目标。当宏观方向由相互冲突的任务共同决定时，简单求和可能抵消成接近零向量，随后归一化放大数值噪声。实现应先检查硬约束，再保留Pareto候选或采用集合值控制；只有在选择规则与量纲换算明确时才形成单一方向。若模型度量从数据学习，还要监控条件数和漂移，防止局部距离因训练分布改变而失真。流形坐标只是内部表示，任何“相近”结论都必须通过任务读出验证。

安全验收要超出范数检查。我们至少测量约束残差、一步角位移、任务风险、校准误差、输入扰动增益、延迟裕度、拒绝率和外部结果；再进行极端输入、重复提示、状态回滚、网络分区和执行失败注入。一个始终位于球面的状态仍可能稳定地产生错误行为，所以外部验收和因果回执不可省略。切空间几何的价值在于明确合法局部方向和减少冗余，它只是闭环系统的一层约束，不是“拓扑防爆”的终极确证。

最后，连续理论与运行监控必须共享同一坐标和时间口径。日志应记录投影前驱动、投影后切向量、回缩结果、步长、度量版本和读出变化，使异常能够定位到编码、几何还是执行接口。只记录最终单位向量会掩盖饱和与巨大抵消，无法解释系统为何转向。对每次高风险更新还应保留未投影候选和拒绝原因，以便离线复核并改进门控。

\section{快慢分离：窗口聚合、随机逼近与漂移跟踪}
\label{sec:macro_adiabatic_separation}

宏观状态通常比局部活动更新得慢，但这种分离应由可测时间尺度和稳定条件建立，而不是借用波恩--奥本海默近似便自动成立。设快状态$h_t$在毫秒至秒级更新，工作状态$w_t$在秒至分钟级更新，宏观状态$z_t$在更长窗口提交。连续模型可写为
\begin{align}
 \epsilon\dot h &= f(h,w,z,o_t)+\nu_h(t),\\
 \dot w &= g(h,w,z,o_t)+\nu_w(t),\\
 \tau_z\dot z &= k(\bar s_t,z,y_t,G_t)+\nu_z(t),
\end{align}
其中$0<\epsilon\ll1$、$\tau_z\gg1$均相对于共同墙钟定义，$o_t$是观测，$y_t$是时变目标，$\bar s_t$是窗口聚合证据。噪声$\nu$表示模型未解释扰动，其协方差单位必须与对应状态导数一致；它不是“认知温度”或设备热噪的同义词。只有当固定$(w,z,o)$时快子系统存在可识别且一致稳定的吸引集，并且目标、输入和结构变化比快收敛足够慢，才能用平均或准稳态近似消去$h$。

宏观更新不应简单平均所有底层活动。设窗口$W_t$内事件$i$具有来源可靠度$\rho_i$、任务影响$\kappa_i$、新颖度$n_i$和不可逆风险$r_i$，可先构造带来源的充分统计$\bar s_t=\operatorname{Agg}\{(s_i,\rho_i,\kappa_i,n_i,r_i)\}$。高频但低影响的重复事件可被降权，短暂却高风险的告警必须直达宏观回路。聚合器同时输出有效样本量、冲突度和置信区间；当证据互相矛盾或来源高度相关时，不允许只凭数量形成高置信。外部记忆可保存原始事件与聚合谱系，宏观状态只保留任务摘要，从而在需要审计时通过来源指针恢复证据，而不是把全部历史压进一个向量。

离散实现更接近双时间尺度随机逼近：
\begin{align}
 h_{k+1}&=h_k+\alpha_k\widehat f_k,\\
 z_{k+1}&=\operatorname{Proj}_{\mathcal Z}\!\left(z_k+\beta_k\widehat k_k\right),
 \qquad \beta_k/\alpha_k\rightarrow0.
\end{align}
经典收敛结论还要求步长和满足适当条件、噪声为可控鞅差、迭代有界、平均场足够光滑且目标近似静止。持续学习系统往往不满足“目标静止”，因此更现实的指标是跟踪误差而非收敛到单一点。若最优宏观状态$z_t^*$以有限速率漂移，可分析$\|z_t-z_t^*\|$的输入到状态界，把目标漂移、结构跳变、观测误差和计算延迟作为输入。不能由时间常数较大推出系统自动拥有战略定力；过慢会错过环境突变，过快则放大噪声，二者需由任务代价与风险共同调节。

AgentOS实例可以用SPC形成任务相关压缩，用TCE调节增益、探索尺度和提交频率，用CCM监控工作状态拥塞，用Boundary控制外部命令与Offloading通道。然而，这是一套工程分工，不是一般理论的定义。其他系统可以用层级贝叶斯过滤器、模型预测控制、黑板架构或多主体协议实现相似快慢功能。无论实现为何，宏观更新都应具有最小驻留时间、紧急旁路、迟滞阈值、版本检查和回滚策略。紧急旁路只允许经过预先授权的安全信号，避免攻击者用“高风险”标签绕过聚合。

时间尺度还与观测窗口相互作用。窗口太短会把偶然波动当成趋势，窗口太长会掩盖制度变化或环境突变；固定窗口也无法同时处理稀疏高风险事件与密集低风险事件。可使用事件触发与周期触发并行的调度方式：周期回路负责稳定统计，事件回路只在证据越过风险门槛时唤醒，并将触发原因写入日志。门槛必须在校准集上估计，并随误报、漏报和资源预算调整，而不是用单一“粘滞度”解释所有延迟。

验证快慢分离时，我们改变输入频谱、突变幅度和延迟，测量快状态恢复时间、宏观状态角速度、告警漏报、误触发、目标跟踪误差与资源账本。还要在目标突然改变、结构模型升级、外部记忆不可用以及长时间偏置输入下重复测试。若宏观变量只是因为过度平滑而显得稳定，却无法追踪任务变化，这不是鲁棒性；若它对每次局部波动都响应，也没有形成真正的层级控制。有效的慢变量是在规定窗口内压制无关波动，同时保留对关键变化的可辨识响应，其成立范围必须随任务和环境一起报告。

\section{自指闭环：固定点、极限环与身份连续性}
\label{sec:strange_loops_and_fixed_points}

宏观状态影响底层处理，底层结果又更新宏观状态，这构成普通的闭环自指，并不需要先假定克莱因瓶或奇异环。设联合闭环状态$q_k=(h_k,G_k,z_k,e_k,m_k)$，控制接口产生$a_k=\pi(z_k,h_k,G_k)$，环境与内部状态经$F$转移，观测经$O$生成，宏观更新经$R$完成，则
\begin{equation}
 q_{k+1}=\mathcal C_k(q_k,\omega_k)
 =\bigl(F_h(q_k,a_k),F_G(q_k,a_k),R(q_k),F_e(e_k,a_k),F_m(q_k)\bigr),
\end{equation}
其中$\omega_k$汇总随机性和未建模输入。这个方程显式保留环境、结构与记忆，不把全部反馈折叠成一个神秘算子。自指的实际困难是模型用自身输出改变后续数据分布、评价标准和可行动作，容易形成确认偏差；解决办法是保留外部参照、反事实测试、独立回执和版本化评价，而不是把闭环形状本身当作同一性的证明。

巴拿赫不动点定理只在完备度量空间中、对同一个压缩映射且压缩常数严格小于一时保证唯一不动点。耗散、阻尼或局部投影都不能单独推出这一条件。若在区域$\mathcal V$内，联合映射各模块的增益矩阵为$K$，并能找到加权范数使$\|K\|<1$，且$\mathcal C(\mathcal V)\subseteq\mathcal V$，才可得到局部唯一固定点及收敛。结构切换、时延、饱和与随机输入会改变$K$，需要公共Lyapunov函数、平均驻留时间或输入到状态分析。更重要的是，固定点与极限环不是同一个结论：压缩映射趋向固定点，非平凡周期轨道则要求研究Poincar\'e映射或分岔条件，不能由巴拿赫定理同时推出。

智能运行也未必应收敛到固定人格点。开放环境中的主体需要在保持承诺的同时学习，身份连续性更适合定义为一组任务相关不变量和允许变化范围。我们令$\mathcal I$包含主体标识、关键承诺、权限边界、来源谱系和长期目标版本；令$\mathcal V_{\mathrm{id}}$为满足这些约束的可行集。只要轨迹在扰动下保持于$\mathcal V_{\mathrm{id}}$，关键查询的变化有界，并且每次不可逆修改都具有授权与审计链，就可以称系统维持了操作意义上的身份连续。周期性行为、稳态分布或缓慢漂移都可能满足这一条件，不必被压缩成唯一极限环。

闭环失败至少有四类。第一，正反馈放大：宏观偏好强化选择性观测，选择性观测又提高该偏好置信。第二，评价自污染：系统修改了用于判断自身成功的度量。第三，结构锁定：重复输出被写入记忆，随后作为独立证据再次读取。第四，环境反作用：控制改变环境后，系统把自身造成的结果误判为外界支持。对应的防护包括保留原始观测、区分自产与外来证据、对指标修改设置更高权限、对记忆回读去重、以及使用对照或反事实估计行动效应。局部结构网络负责保存因果与版本谱系，全局分布表示不能以高置信覆盖来源冲突。

对同一性还要区分技术连续与伦理连续。复制一个状态快照可能保持大量查询，却未必继承原主体的权限、责任或社会承诺；合并两个记忆分支也可能产生无法调和的版本冲突。因此，身份不是向量距离的单一阈值，而是由技术标识、行为连续、授权链和制度承认共同构成。一般状态动力学可以描述这些变量如何变化，却不能单独决定法律或伦理上的主体资格，相关判断必须由外部制度接口给出。

验证时，我们分别寻找不动点、周期轨道、可行集保持和漂移界，不把一种稳定概念替代另一种。测试应包括共同输入遮蔽、评价函数切换、记忆回灌、网络分区、控制饱和和环境延迟；并用外部任务结果、规范合规与恢复能力三套指标验收。内部损失下降只能说明指定目标在指定数据上的变化，不保证事实正确、社会可接受或任务最终完成。由此，“流体自我”可以保留为持续更新而保持约束的解释性名称，其数学对象是版本化闭环轨迹与身份可行集，而不是由物理耗散必然产生的永恒几何灵魂。

\section{向下控制：宏观策略接口、探索调节与外部验收}
\label{sec:output_characteristics_gauge_generator}

宏观状态的价值最终要通过对底层计算与外部行动的可测影响体现。我们把所谓“向下因果”定义为跨层接口上的条件控制，而不是宏观变量修改底层物理定律。设$z_t$经策略映射$B_{\phi}$产生控制包
\begin{equation}
 c_t=B_{\phi}(z_t,\Sigma_t,\mathcal R_t)
 =\bigl(g_t^{\mathrm{att}},g_t^{\mathrm{inh}},\tau_t^{\mathrm{exp}},b_t^{\mathrm{route}},\Gamma_t^{\mathrm{act}}\bigr),
\end{equation}
其中$\Sigma_t$是宏观不确定性，$\mathcal R_t$是资源与权限账本；五个分量分别表示注意增益、抑制增益、探索尺度、路由偏置和允许动作集合。它们都有明确定义域与饱和范围，通过Boundary进入下层算子。控制包改变的是候选权重、采样分布、计算配额、工具路由和动作资格，不是哈密顿量、重力或热力学温度。若某个具身系统确实调节物理执行器或设备功率，应另以工程单位建模并通过传感器计量。

以全局分布读出为例，候选$j$的logit可写为
\begin{equation}
 \ell_{j,t}'=\operatorname{clip}\!\left(\ell_{j,t}
 +g_t^{\mathrm{att}}a_{j,t}-g_t^{\mathrm{inh}}r_{j,t}
 +b_{j,t}^{\mathrm{route}},\,-L,L\right),
 \qquad p_{j,t}=\operatorname{softmax}(\ell_t'/\tau_t^{\mathrm{exp}})_j.
\end{equation}
这里$a_{j,t}$和$r_{j,t}$是有来源的支持与风险证据，$L$限制数值范围，$\tau^{\mathrm{exp}}>0$是无量纲采样尺度。提高探索尺度通常使分布更平坦，但熵的具体变化依候选和截断而定；它不等于向系统注入热能。降低尺度也不会把灵感“凝固”进结构，结构写回必须经过重复证据、保留查询、冲突检测、授权和回滚。注意、抑制、路由与结构学习是不同接口，不能用一个温度旋钮替代全部机制。

在AgentOS实例中，SPC可从$h(t)$、$E$、$\Theta$及外部记忆$M_e$提取任务相关宏观摘要，TCE根据该摘要调节探索、增益和控制力度，CCM监控候选分支、循环、资源与延迟，Boundary限制命令的作用范围，Offloading把高代价状态迁移到工具、文件、环境或其他Agent。$\Psi$可作为主体相关的优先级与长期承诺状态参与映射，但不能凌驾于事实证据和安全约束。局部结构网络提供“控制针对哪个对象、关系和版本”，全局分布表示提供“当前候选及其不确定性”；只有绑定成功的控制才允许提交，绑定缺失时应拒绝而非广播模糊偏置。

跨层控制必须处理权限和责任。宏观层无权直接覆盖所有局部约束：底层执行器应校验类型、单位、范围、租约和前置条件；高风险动作需要二次确认或外部授权；社会系统中的下层主体还保有拒绝、申诉和退出接口。层级控制的优势来自摘要与协调，不意味着高层状态天然更真实、更有能力或拥有无限权力。若宏观摘要遗漏局部关键变量，底层必须上报残差；若命令不可行，应返回结构化失败而非静默执行。这样，控制层级才是可纠错协议，而非“创世算子”。

控制策略还应区分事实不确定性、价值冲突与资源不足。事实不确定时应请求观测或降低动作可逆性；价值冲突时应保留多方案并调用规范仲裁；资源不足时则调整期限、卸载或拒绝任务。把三者都解释为“升温探索”会使系统在需要谨慎核实时盲目增加随机性。宏观接口应输出原因码与证据摘要，使下层和外部监督者知道控制变化源于何种诊断，并能检查诊断是否成立。

因果效果要靠干预与回执识别。记录执行前状态、控制包、实际执行量、环境响应和独立验收结果，并与无控制、替代控制或延迟控制比较，估计任务风险差。动作发出不等于动作发生，环境改变也不等于系统理解了环境；结果可能来自共同输入、其他主体或测量偏差。任务成功、规范合规、资源消耗和设备能耗应分别报告。预测有效不能推出整体同构，内部目标下降不能推出外部成功，控制影响也不能推出价值正当。

本章最终得到的是一条可审计链：高维联合状态在任务查询下形成宏观压缩，宏观变量依据适用条件选择几何，历史更新通过连续或离散算子积累，约束动力学在明确步长下实现，快慢回路以跟踪误差与风险调节，闭环轨迹在身份可行集内保持，控制包再经权限接口作用于底层和环境。每一环都可以替换、反驳和实验，不依赖超球面、李群或物理隐喻成为普遍真理。HSF-HD由此维持“智能的一般空气动力学”定位：它研究状态、结构、目标、资源与环境如何共同塑造轨迹；具体AgentOS则把这些原则落实为SPC、TCE、CCM、Boundary、Offloading和外部记忆等工程组件。

%---chp---
